\documentclass[11pt, a4paper]{amsart}
\usepackage{amsmath, amssymb, amsthm, amsfonts}
\usepackage{mathtools}
\usepackage[margin=1in]{geometry}
\usepackage{hyperref}
\usepackage{tikz-cd}
\usepackage{enumitem}

\newtheorem{theorem}{Theorem}[section]
\newtheorem{lemma}[theorem]{Lemma}
\newtheorem{proposition}[theorem]{Proposition}
\newtheorem{corollary}[theorem]{Corollary}
\newtheorem{fact}[theorem]{Standard Fact}
\theoremstyle{definition}
\newtheorem{definition}[theorem]{Definition}
\newtheorem{remark}[theorem]{Remark}
\newtheorem{notation}[theorem]{Notation}
\newtheorem{example}[theorem]{Example}

\newcommand{\R}{\mathbb{R}}
\newcommand{\Q}{\mathbb{Q}}
\newcommand{\C}{\mathbb{C}}
\newcommand{\Z}{\mathbb{Z}}
\newcommand{\N}{\mathbb{N}}
\newcommand{\HH}{\mathbb{H}}
\newcommand{\OO}{\mathbb{O}}
\newcommand{\id}{\operatorname{id}}
\newcommand{\im}{\operatorname{im}}
\newcommand{\tr}{\operatorname{tr}}
\newcommand{\Cl}{\mathrm{Cl}}
\newcommand{\inv}{\mathrm{inv}}
\newcommand{\cyl}{\mathrm{cyl}}
\renewcommand{\Im}{\mathrm{Im}}

\title{Cascade--Classical Correspondence III:\\
  Refinement Spaces and Symmetry Actions}
\author{}
\date{}

\begin{document}

\maketitle

\begin{abstract}
This paper builds the finite-level analytic infrastructure
intended for downstream cascade-correspondence papers. From the
algebraic substrate of paper P2, we construct: iterated refinement
data $\mathcal{R}_n^\bullet = (V_n^\bullet, E_n^\bullet, F_n^\bullet,
\partial_n^\bullet)$ for $\bullet \in \{H_4, D_4\}$, including
explicit $2$-face data needed to place centroid vertices, with
literal vertex inclusion $V(G_n^\bullet) \subset V(G_{n+1}^\bullet)$;
finite-level $\ell^2$ vertex spaces $X_n^0$ and edge $\ell^2$
spaces $X_n^1$ under counting measure; downward bonding maps $p_{n+1, n}^0$ defined
as restriction to $V(G_n^\bullet)$ and $p_{n+1, n}^1$ defined as
fibrewise averaging over the canonical edge-parent map; the
projective inverse-limit Hilbert space (in the vertex sector)
$X_\inv^{0, \bullet} \cong \ell^2(V_\infty^\bullet; F^0)$; the
cylindrical inductive-limit space $X_\cyl^{0, \bullet} =
c_{00}(V_\infty^\bullet; F^0)$ embedded densely in
$X_\inv^{0, \bullet}$; full intertwining identities for the
refinement-Coxeter and refinement-$\sigma$ pairs; and a restricted
intertwining identity for the Coxeter-$\sigma$ pair (full
$W(H_4)$-$\sigma$ commutation on the $V_{H_4}$ sector fails because
$W(H_4)$ acts via $\Q(\varphi)$-linear matrices; see
Theorem~\ref{thm:commute}(4) and
Remark~\ref{rmk:coxeter-sigma-restricted} for the precise scope).

This is the third of thirteen planned cascade infrastructure papers.
The construction is in the correct projective direction: bonding
maps $p_{n+1, n}$ go \emph{downward} from level $n+1$ to level $n$,
implemented as restriction (vertex sector) or as canonical
parent-fibre averaging (edge sector). Their adjoints
$j_{n, n+1} := p_{n+1, n}^*$ are (rescaled) embeddings going upward.
In the vertex sector the adjoint is a genuine section
($p^0 j^0 = \id$), the cylindrical direct limit is well-defined,
and the inverse limit is its Hilbert completion. In the edge sector
the adjoint is only a \emph{half-section} ($p^1 j^1 = \tfrac12 \id$),
a clean cylindrical direct-limit / Hilbert-completion package is
not established here, and edge-sector statements are kept at the
finite level with bounded bonding maps only
(Remarks~\ref{rmk:edge-cyl} and~\ref{rmk:edge-inv}). \emph{No continuum-limit, Gromov--Hausdorff,
spectral-convergence, or manifold-identification claim is made
here.}

The main theorems proved in this paper:
(1) Proposition~\ref{prop:refinement-equivariant}, the refinement
rule with $V_n \subset V_{n+1}$ inclusion is well-defined and
Coxeter-equivariant;
(2) Proposition~\ref{prop:Xn-Hilbert}, each finite-level $X_n^0,
X_n^1$ is a finite-dimensional Hilbert space;
(3) Theorem~\ref{thm:bonding}, downward bonding maps are
contractions under counting measure;
(4) Proposition~\ref{prop:adjoints}, the cylindrical embeddings
$j_{n, n+1}$ are the adjoints of the bonding maps;
(5) Theorem~\ref{thm:Xcyl-direct-limit}, the vertex algebraic
direct limit $X_\cyl^{0, \bullet}$ is well-defined as
$c_{00}(V_\infty^\bullet; F^0)$;
(6) Theorem~\ref{thm:Xinv-Hilbert}, the vertex inverse limit
$X_\inv^{0, \bullet}$ under restriction is canonically isometric
to $\ell^2(V_\infty^\bullet; F^0)$ and is a separable Hilbert space;
(7) Proposition~\ref{prop:Xcyl-dense}, $X_\cyl^{0, \bullet}$
embeds densely in $X_\inv^{0, \bullet}$;
(8) Theorem~\ref{thm:commute}, the commutation
$g \circ p_{n+1, n} = p_{n+1, n} \circ g$ for $g$ in the relevant
Coxeter group, and analogously for $\sigma$ from P1.

These are the load-bearing inputs for the closure-dynamics paper P4,
the metric projection paper P5, and the gauge/arithmetic/computation
branch papers P9--P13, as well as for the correspondence Foundations
paper.
\end{abstract}

\tableofcontents

%=====================================================================
\section{Introduction and Source Audit}
\label{sec:intro}
%=====================================================================

\subsection{The repair target}

An earlier draft of the correspondence Foundations manuscript
defined the cascade limit Hilbert space using forward zero-extension
maps $X_n \to X_{n+1}$ and called the resulting object an
``inverse limit''. That construction is incorrect on two counts:
the system is inductive (not projective), and forward zero-extension
is not isometric under any naturally-normalized $\ell^2$ structure on
finite refinement levels.

The correct construction is built on \emph{downward} bonding maps
$p_{n+1, n} \colon X_{n+1} \to X_n$. The cleanest mathematical
choice for these bonding maps, used here, is:
\begin{itemize}
\item Vertex sector: $p_{n+1, n}^0$ is restriction of a function on
  $V(G_{n+1}^\bullet)$ to its sub-domain $V(G_n^\bullet)
  \subset V(G_{n+1}^\bullet)$.
\item Edge sector: $p_{n+1, n}^1$ is fibrewise averaging along the
  canonical edge-parent map $E^{\rm or}(G_{n+1}^\bullet) \to
  E^{\rm or}(G_n^\bullet)$, which assigns each subdivided
  level-$(n+1)$ oriented edge to the level-$n$ oriented edge it is
  a piece of.
\end{itemize}

Both maps are genuine contractions under \emph{counting measure}
on finite vertex/edge sets, and their adjoints are the natural
embeddings in the upward direction. The projective system in the
correct direction is then $X_\inv := \{(f_n) : p_{n+1, n} f_{n+1} =
f_n\}$, equipped with the limit inner product, and we show this is
canonically isometric to $\ell^2(V_\infty^\bullet)$ for the
countable union $V_\infty^\bullet := \bigcup_n V(G_n^\bullet)$.

\subsection{What this paper proves}

The substrate constructed here consists of:
\begin{enumerate}
\item Iterated refinement graphs $G_n^{H_4}$ and $G_n^{D_4}$ with
  $V(G_n^\bullet) \subset V(G_{n+1}^\bullet)$ and a canonical
  edge-parent map (\S\ref{sec:graphs}).
\item Counting measures on vertices and oriented edges
  (\S\ref{sec:measures}).
\item Finite-level vertex and edge $\ell^2$ spaces $X_n^0, X_n^1$,
  each a finite-dimensional Hilbert space
  (\S\ref{sec:finite-level}).
\item Downward bonding maps $p_{n+1, n}^0$ (restriction) and
  $p_{n+1, n}^1$ (parent-fibre averaging), both contractions
  (\S\ref{sec:bonding}).
\item Cylindrical embeddings $j_{n, n+1}^0$ (zero-extension) and
  $j_{n, n+1}^1$ (rescaled replication) as adjoints of the bonding
  maps (\S\ref{sec:bonding}).
\item Algebraic direct-limit space (vertex sector)
  $X_\cyl^{0, \bullet} \cong c_{00}(V_\infty^\bullet; F^0)$
  (\S\ref{sec:cyl-lim}).
\item Projective inverse-limit Hilbert space (vertex sector)
  $X_\inv^{0, \bullet} \cong \ell^2(V_\infty^\bullet; F^0)$
  (\S\ref{sec:inv-lim}).
\item Commuting refinement / Coxeter / $\sigma$ intertwining
  identities for $p_{n+1, n}$ and $j_{n, n+1}$
  (\S\ref{sec:commute}).
\end{enumerate}

\subsection{What this paper does not prove}

P3 is deliberately narrow. We record the exclusions explicitly:
\begin{enumerate}
\item No continuum-limit, Gromov--Hausdorff, or spectral-convergence
  statement.
\item No closure-functional, gradient-flow, energy, or
  monotonicity theorem.
\item No identification of $X_\inv$ with any classical continuum
  function space such as $L^2(S^3)$.
\item No metric, hydrodynamic, gauge, spectral, zeta, Hecke, or
  CTM/TM correspondence theorem.
\item No vertex-wise quotient map $V(G_n^{H_4}) \to V(G_n^{D_4})$.
\item No physical or phenomenological interpretation.
\end{enumerate}

\subsection{Source audit}

Theorem inputs to this paper are:
\begin{description}
\item[\textbf{Internal cascade infrastructure.}]
Paper P1 (cascade--classical correspondence I, $\sigma$-fixed
rationality \cite{SigmaRationality}) for the coefficientwise
$\sigma$-action; paper P2 (cascade algebraic substrate
\cite{AlgebraicSubstrate}) for the algebraic substrate, including
$D_4 \subset E_8$, the icosian model of $H_4$, the explicit
$V_{600}$ and $V_{24}$ vertex sets, and the $\Cl(1, 3)$ and
octonionic sector models.
\item[\textbf{Classical mathematics.}]
Standard Hilbert space and direct/inverse limit theory
\cite{ReedSimonI, Bogachev}; standard Coxeter group theory
\cite{Humphreys, CoxeterRegularPolytopes}; finite graph and
$1$-cochain conventions \cite{Hatcher}; elementary $\ell^2$ theory.
\end{description}

We do not cite as theorem-grade input any other cascade-internal
note. The bibliography below distinguishes load-bearing
upstream sources (P1, P2) used as theorem-grade input from
downstream cascade preprints
(\cite{CascadeClosureDynamics, CascadeRefinementCompat,
CascadeMechanism}) that consume the present paper's results, and
classical references used as background or auxiliary input.

\subsection{Organization}

Section~\ref{sec:graphs} sets up the refinement graphs and the
canonical edge-parent map. Section~\ref{sec:measures} fixes the
counting measures. Section~\ref{sec:finite-level} defines the
finite-level vertex and edge $\ell^2$ spaces.
Section~\ref{sec:bonding} constructs the downward bonding maps and
their adjoints, the cylindrical embeddings.
Section~\ref{sec:cyl-lim} treats the cylindrical algebraic
direct-limit space. Section~\ref{sec:inv-lim} constructs the
projective limit Hilbert space and proves it is $\ell^2(V_\infty)$.
Section~\ref{sec:commute} establishes the intertwining identities
for refinement, Coxeter, and $\sigma$. Section~\ref{sec:handoff}
states the citation contract for downstream papers.
Appendix~\ref{app:notation} is a notation concordance.

%=====================================================================
\section{Refinement Graphs and the Edge-Parent Map}
\label{sec:graphs}
%=====================================================================

\subsection{The base graphs}

\begin{definition}[Base graphs]
\label{def:base-graphs}
Let $V_{600} \subset \HH \cong \R^4$ be the standard $120$-vertex
set of the regular $600$-cell, in the coordinate realization of
\cite[\S 5]{AlgebraicSubstrate}. Let $V_{24} \subset V_{600}$ be
the $24$-vertex subset (axis points and half-type points) congruent
to the rescaled $D_4$ root polytope by
\cite[Theorem~\texttt{thm:shell-class}]{AlgebraicSubstrate}.

The \emph{base $H_4$ graph} $G_0^{H_4} = (V_0^{H_4}, E_0^{H_4})$
has vertex set $V_0^{H_4} := V_{600}$ and unoriented edges joining
pairs of $600$-cell vertices at the standard $600$-cell edge
distance \cite[\S 22.7]{CoxeterRegularPolytopes}.

The \emph{base $D_4$ graph} $G_0^{D_4} = (V_0^{D_4}, E_0^{D_4})$
has vertex set $V_0^{D_4} := V_{24}$ and unoriented edges joining
pairs of $24$-cell vertices at the standard $24$-cell edge
distance \cite[\S 22.2]{CoxeterRegularPolytopes}.

We write $\bullet$ for either decoration and $G_n^\bullet$ for the
level-$n$ graph below. The Coxeter group $W^\bullet$ is $W(H_4)$
when $\bullet = H_4$ and $W(D_4)$ when $\bullet = D_4$.
\end{definition}

\subsection{Refinement with vertex inclusion}

A graph alone does not determine the Schl\"{a}fli refinement; the
construction needs explicit $2$-face data, since centroids are
indexed by $2$-faces. We carry that data through levels as a
\emph{refinement datum}.

\begin{definition}[Refinement datum]
\label{def:refinement-datum}
A \emph{refinement datum at level $n$} is a tuple
$\mathcal{R}_n^\bullet = (V_n^\bullet, E_n^\bullet, F_n^\bullet,
\partial_n^\bullet)$ where $V_n^\bullet$ and $E_n^\bullet$ are
the vertex and unoriented edge sets of $G_n^\bullet$;
$F_n^\bullet$ is a finite set of \emph{$2$-faces} (each face is
a finite cyclic sequence of edges
$(e_1, \ldots, e_k) \in E_n^\bullet$ with $k \ge 3$ closing up
into a polygon, with the standard cyclic-equivalence convention);
and $\partial_n^\bullet \colon F_n^\bullet \to \bigsqcup_k
(E_n^\bullet)^k$ is the boundary map sending each face to its
edge sequence. The base data $\mathcal{R}_0^\bullet$ is induced
by Definition~\ref{def:base-graphs} from the $1$-skeleton and
$2$-faces of the underlying $4$-polytope ($V_{600}$ for
$\bullet = H_4$, $V_{24}$ for $\bullet = D_4$). $W^\bullet$ acts
on each component of $\mathcal{R}_n^\bullet$ compatibly with the
boundary map.
\end{definition}

\begin{definition}[Schl\"{a}fli refinement with vertex inclusion]
\label{def:refinement}
Given a refinement datum
$\mathcal{R}_n^\bullet = (V_n^\bullet, E_n^\bullet, F_n^\bullet,
\partial_n^\bullet)$, the \emph{Schl\"{a}fli refinement} produces
a new datum $\mathcal{R}_{n+1}^\bullet$ with the following
explicit components:
\begin{itemize}
\item \textbf{Vertices.}
  $V_{n+1}^\bullet := V_n^\bullet \sqcup M_n^\bullet \sqcup
  C_n^\bullet$, where $M_n^\bullet$ contains exactly one
  edge-midpoint $m_e$ per $e \in E_n^\bullet$, and $C_n^\bullet$
  contains exactly one face-centroid $c_f$ per $f \in F_n^\bullet$.
  The first summand is the literal inclusion
  $V_n^\bullet \subset V_{n+1}^\bullet$.
\item \textbf{Edges.} $E_{n+1}^\bullet$ is the disjoint union of
  three explicit edge classes:
  \begin{enumerate}[label=(\alph*), leftmargin=2.4em, nosep, topsep=1pt]
  \item \emph{Subdividing edges:} for each $e = \{u, v\} \in
    E_n^\bullet$, the two edges $\{u, m_e\}$ and $\{m_e, v\}$.
  \item \emph{Midpoint-centroid edges:} for each face
    $f \in F_n^\bullet$ with edge sequence
    $(e_1, \ldots, e_k) = \partial_n^\bullet(f)$, the $k$ edges
    $\{m_{e_i}, c_f\}$ for $i = 1, \ldots, k$.
  \item \emph{Midpoint-midpoint diagonal edges:} for each face
    $f \in F_n^\bullet$ with edge sequence $(e_1, \ldots, e_k)$,
    the $k$ edges $\{m_{e_i}, m_{e_{i+1}}\}$ (cyclic indices),
    closing each level-$(n+1)$ triangle face below.
  \end{enumerate}
\item \textbf{Faces.} $F_{n+1}^\bullet$ is the disjoint union of
  level-$(n+1)$ triangles: each level-$n$ face $f$ with edge
  sequence $(e_1, \ldots, e_k)$ produces $k$ triangles
  $(m_{e_i}, c_f, m_{e_{i+1}})$ (cyclic indices), each with edge
  sequence
  $(\{m_{e_i}, c_f\}, \{c_f, m_{e_{i+1}}\}, \{m_{e_{i+1}}, m_{e_i}\})$
  drawn from edge classes (b), (b), (c) respectively. All three
  boundary edges lie in $E_{n+1}^\bullet$ by construction, so
  $\partial_{n+1}^\bullet$ is well-defined and the iterated
  refinement datum is well-formed.
\item \textbf{$W^\bullet$-action.} $W^\bullet$ acts on
  $V_{n+1}^\bullet$ by $g \cdot v = gv$ for $v \in V_n^\bullet$,
  $g \cdot m_e = m_{ge}$ for $m_e \in M_n^\bullet$ (well-defined
  since $g\{u, v\} = \{gu, gv\}$), and $g \cdot c_f = c_{gf}$
  for $c_f \in C_n^\bullet$ (well-defined since $W^\bullet$ acts
  on $F_n^\bullet$). The induced action on $E_{n+1}^\bullet$
  permutes each of the three edge classes (a), (b), (c) among
  themselves and on $F_{n+1}^\bullet$ permutes triangles among
  themselves.
\end{itemize}
\end{definition}

\begin{remark}[Scope of the construction]
\label{rmk:refinement-scope}
The bonding maps, chain-space, and intertwining constructions
below depend only on the level-$(n+1)$ vertex and edge data, and in
particular only on edge classes (a) and (b) (the subdividing and
midpoint-centroid edges, since the edge sector
$X_n^{1, \bullet}$ averaging is restricted to subdividing edges
via Definition~\ref{def:edge-parent}). Edge class (c)
(midpoint-midpoint diagonals) and the level-$(n+1)$ face
structure $F_{n+1}^\bullet$ are present in the refinement datum
to keep iteration well-formed but are not load-bearing for the
intertwining theorem (Theorem~\ref{thm:commute}). Downstream
papers \cite{CascadeClosureDynamics, CascadeRefinementCompat}
either inherit the full datum or use a deliberately reduced
``pure-midpoint'' abstract model that omits the face structure;
\cite{CascadeRefinementCompat}~\S1.5 names the gap explicitly.
\end{remark}

\begin{remark}[Vertex inclusion is the key]
\label{rmk:vertex-inclusion}
The literal subset $V_n^\bullet \subset V_{n+1}^\bullet$ from
Definition~\ref{def:refinement} is the structural feature that
makes the entire bonding-map construction below work cleanly:
restriction of a function on $V_{n+1}^\bullet$ to $V_n^\bullet$ is
a well-defined function on $V_n^\bullet$, and (under counting
measure) is automatically a contraction in the $\ell^2$-norm.
\end{remark}

\subsection{The canonical edge-parent map}

Unlike vertices, edges are not preserved under refinement: each
level-$n$ edge $\{u, v\} \in E_n^\bullet$ is replaced by two
level-$(n+1)$ edges $\{u, m\}$ and $\{m, v\}$. We record the
inverse direction as a function:

\begin{definition}[Edge-parent map]
\label{def:edge-parent}
For oriented edges, define
\[
  \pi_{n+1, n}^E \colon E^{\rm or}(G_{n+1}^\bullet) \to
    E^{\rm or}(G_n^\bullet) \cup \{\bot\}
\]
by:
\begin{itemize}
\item If $e' \in E^{\rm or}(G_{n+1}^\bullet)$ is one of the two
  pieces of a subdivided level-$n$ oriented edge $e =
  (u, v) \in E^{\rm or}(G_n^\bullet)$ --- i.e., $e' = (u, m)$ or
  $e' = (m, v)$ where $m$ is the midpoint of $\{u, v\}$ ---
  then $\pi_{n+1, n}^E(e') := e$ (with the orientation inherited).
\item If $e'$ is one of the new level-$(n+1)$ edges connecting a
  midpoint to a centroid or two midpoints (i.e., a ``boundary''
  edge introduced by Schl\"{a}fli refinement that does not
  subdivide a level-$n$ edge), then $\pi_{n+1, n}^E(e') := \bot$.
\end{itemize}
The edge-parent map is well-defined as a function and is
$W^\bullet$-equivariant: $g \cdot \pi_{n+1, n}^E(e') = \pi_{n+1, n}^
E(g \cdot e')$ for every $g \in W^\bullet$ and every $e' \in
E^{\rm or}(G_{n+1}^\bullet)$. This holds because Schl\"{a}fli
refinement is itself $W^\bullet$-equivariant, so $W^\bullet$ permutes
the descendants of an edge among themselves.
\end{definition}

\begin{notation}[Subdividing edges]
\label{not:subdividing}
We write $E^{\rm or}_{\rm sub}(G_{n+1}^\bullet) := \{e' :
\pi_{n+1, n}^E(e') \neq \bot\}$ for the set of level-$(n+1)$
oriented edges that subdivide a level-$n$ edge. By construction,
$|E^{\rm or}_{\rm sub}(G_{n+1}^\bullet)| = 2 \cdot |E^{\rm or}(G_n^
\bullet)|$ (each level-$n$ edge contributes exactly two
level-$(n+1)$ subdividing edges).
\end{notation}

\subsection{Properties of the refinement}

\begin{proposition}[Refinement properties]
\label{prop:refinement-equivariant}
For every $n \geq 0$ and $\bullet \in \{H_4, D_4\}$:
\begin{enumerate}
\item $V_n^\bullet \subset V_{n+1}^\bullet$ is a literal subset
  inclusion of finite sets.
\item The action of $W^\bullet$ on $V_n^\bullet$ extends to an
  action on $V_{n+1}^\bullet$ that commutes with the inclusion.
\item The edge-parent map $\pi_{n+1, n}^E$ is well-defined and
  $W^\bullet$-equivariant.
\item Each level-$n$ oriented edge has exactly two level-$(n+1)$
  preimages under $\pi_{n+1, n}^E$ (its two subdividing pieces in
  the inherited orientation).
\end{enumerate}
\end{proposition}

\begin{proof}
(1) and (2) follow directly from Definition~\ref{def:refinement}:
the vertex inclusion is by construction, and the
$W^\bullet$-action extends because midpoints and centroids are
$W^\bullet$-equivariantly defined from the level-$n$ data.

(3) Edge-parent equivariance: by the Schl\"{a}fli refinement
construction \cite[\S 8.2]{CoxeterRegularPolytopes}, $W^\bullet$
permutes subdivided edges among themselves, and the
parent-edge assignment respects this permutation.

(4) Each level-$n$ unoriented edge $\{u, v\}$ contributes
oriented edges $(u, v)$ and $(v, u)$ at level $n$, each of which
has two subdividing pieces at level $n+1$ (namely $(u, m), (m, v)$
for $(u, v)$, and $(v, m), (m, u)$ for $(v, u)$). So there are
two preimages per oriented edge, giving the $|E^{\rm or}_{\rm sub}
(G_{n+1}^\bullet)| = 2 |E^{\rm or}(G_n^\bullet)|$ count of
Notation~\ref{not:subdividing}.
\end{proof}

\begin{remark}[No non-canonical choices]
\label{rmk:no-canonical-choice}
Definition~\ref{def:edge-parent} introduces no non-canonical
choice: each level-$(n+1)$ oriented edge is either intrinsically a
subdividing piece of a unique level-$n$ oriented edge (in which
case it has a unique parent), or it is intrinsically a boundary
edge of the Schl\"{a}fli refinement (in which case it has parent
$\bot$). No lex-order, evaluation, or vertex-choice rule is needed.
\end{remark}

%=====================================================================
\section{Counting Measures}
\label{sec:measures}
%=====================================================================

\begin{definition}[Counting measures]
\label{def:counting-measures}
For each $n \geq 0$ and $\bullet \in \{H_4, D_4\}$:
\begin{align*}
  \mu_n^\bullet(S) &:= |S|, \qquad S \subseteq V(G_n^\bullet); \\
  \nu_n^\bullet(T) &:= |T|, \qquad T \subseteq E^{\rm or}(G_n^
    \bullet).
\end{align*}
These are unweighted counting measures; in particular, every vertex
and oriented edge has measure $1$.
\end{definition}

\begin{remark}[Why counting measure]
\label{rmk:why-counting}
We use unscaled counting measure rather than probability
normalization. The reason is that, with $V_n^\bullet \subset
V_{n+1}^\bullet$, restriction to $V_n^\bullet$ is a contraction
in $\ell^2(V_{n+1}^\bullet, \mu_{n+1}^\bullet) \to
\ell^2(V_n^\bullet, \mu_n^\bullet)$ \emph{automatically}: the
norm of the restriction is $\sum_{v \in V_n^\bullet} |f(v)|^2$,
which is at most $\sum_{w \in V_{n+1}^\bullet} |f(w)|^2 =
\|f\|^2_{X_{n+1}^0}$ since the smaller sum has fewer terms (all
nonneg). Probability normalization would introduce a level-dependent
rescaling factor that breaks this property.
\end{remark}

\begin{notation}[Sector fibres]
\label{not:fibres}
We import from \cite[\S\S 3, 5, 6, 7]{AlgebraicSubstrate}:
\begin{itemize}
\item $V_{H_4} \cong \HH \cong \R^4$ with the standard Euclidean
  inner product (square-norm of unit icosian $= 1$);
\item $V_{D_4} \cong \R^4$ with the standard Euclidean inner
  product;
\item $V_{16} = \Cl(1, 3)$ with the Hilbert--Schmidt-style trace
  form $\langle a, b \rangle_{\Cl} := \langle a^\dagger b
  \rangle_0$ (positive-definite by \cite[Ch.~9]{Lounesto});
\item $V_8 = \OO$ with the standard octonion norm form
  $\langle a, b \rangle_\OO := \Re(a \bar b)$.
\end{itemize}
We fix the combined vertex fibre
\[
  F^0 := V_{H_4} \oplus V_{D_4} \oplus V_{16} \oplus V_8,
\]
$\R$-dimension $4 + 4 + 16 + 8 = 32$, with the orthogonal-direct-sum
inner product, and the edge fibre
\[
  F^1 := \Im(\OO),
\]
$\R$-dimension $7$, with the Euclidean inner product induced from
$\OO$. These choices are fixed throughout this paper. (Downstream
papers needing a different edge fibre re-construct the relevant
spaces from scratch; we make no template promise to accommodate
arbitrary $F^1$.)
\end{notation}

%=====================================================================
\section{Sector Fibres and Finite-Level $\ell^2$ Spaces}
\label{sec:finite-level}
%=====================================================================

\begin{definition}[Vertex-sector $\ell^2$ space]
\label{def:Xn-0}
For $n \geq 0$ and $\bullet \in \{H_4, D_4\}$, define
\[
  X_n^{0, \bullet} := \ell^2\bigl(V(G_n^\bullet), \mu_n^\bullet;\,
    F^0\bigr),
\]
the $F^0$-valued $\ell^2$-functions on $V(G_n^\bullet)$ under
counting measure, with inner product
\[
  \langle f, g \rangle_{X_n^0} :=
    \sum_{v \in V(G_n^\bullet)} \langle f(v), g(v) \rangle_{F^0}.
\]
\end{definition}

\begin{definition}[Edge-sector $\ell^2$ space]
\label{def:Xn-1}
For $n \geq 0$ define
\[
  X_n^{1, \bullet} := \ell^2_{\rm anti}\bigl(E^{\rm or}(G_n^\bullet),
    \nu_n^\bullet;\, F^1\bigr),
\]
the $F^1$-valued $\ell^2$-functions $A$ on oriented edges
satisfying the antisymmetry convention $A(-e) = -A(e)$ for every
oriented edge $e$ (where $-e$ denotes reversed orientation), with
inner product
\[
  \langle A, B \rangle_{X_n^1} := \tfrac{1}{2}
    \sum_{e \in E^{\rm or}(G_n^\bullet)} \langle A(e), B(e)
    \rangle_{F^1}.
\]
The factor of $\tfrac{1}{2}$ accounts for the double-counting from
including both orientations of each unoriented edge.
\end{definition}

\begin{proposition}[Finite-level Hilbert structure]
\label{prop:Xn-Hilbert}
For every $n \geq 0$ and $\bullet \in \{H_4, D_4\}$:
\begin{enumerate}
\item $X_n^{0, \bullet}$ and $X_n^{1, \bullet}$ are
  finite-dimensional real Hilbert spaces.
\item $\dim_\R X_n^{0, \bullet} = |V(G_n^\bullet)| \cdot 32$ and
  $\dim_\R X_n^{1, \bullet} = |E(G_n^\bullet)| \cdot 7$, where
  $|E(G_n^\bullet)|$ is the number of unoriented edges
  ($= |E^{\rm or}(G_n^\bullet)|/2$).
\item In particular, both spaces are separable Hilbert spaces.
\end{enumerate}
\end{proposition}

\begin{proof}
Each fibre is a finite-dimensional Euclidean space (32 for $F^0$,
$7$ for $F^1$), and $V(G_n^\bullet), E^{\rm or}(G_n^\bullet)$ are
finite (Proposition~\ref{prop:refinement-equivariant}). A finite
direct sum of finite-dimensional Euclidean spaces under a
positive-definite bilinear form is a finite-dimensional Hilbert
space \cite[Ch.~II]{ReedSimonI}. The dimension formulas follow by
counting; the antisymmetry convention identifies each pair
$(e, -e)$ with a single element of dimension $\dim_\R F^1 = 7$,
giving the unoriented-edge count.
Separability is automatic.
\end{proof}

%=====================================================================
\section{Downward Bonding Maps and Their Adjoints}
\label{sec:bonding}
%=====================================================================

\subsection{Vertex-sector bonding by restriction}

\begin{definition}[Vertex bonding map]
\label{def:p0}
For $n \geq 0$, define the \emph{vertex bonding map}
\[
  p_{n+1, n}^0 \colon X_{n+1}^0 \to X_n^0, \qquad
    (p_{n+1, n}^0 f)(v) := f(v) \quad \text{for } v \in V(G_n^
    \bullet).
\]
Since $V(G_n^\bullet) \subset V(G_{n+1}^\bullet)$ (Definition
\ref{def:refinement}), the right-hand side is well-defined: $f$ is
already a function on $V(G_{n+1}^\bullet)$ and we simply restrict
its domain to the subset $V(G_n^\bullet)$.
\end{definition}

\subsection{Edge-sector bonding by parent-fibre averaging}

\begin{definition}[Edge bonding map]
\label{def:p1}
For $n \geq 0$, define the \emph{edge bonding map}
\[
  p_{n+1, n}^1 \colon X_{n+1}^1 \to X_n^1, \qquad
    (p_{n+1, n}^1 A)(e) := \tfrac{1}{2} \sum_{e' \in (\pi_{n+1, n}^E)^{-1}(e)}
    A(e')
\]
for $e \in E^{\rm or}(G_n^\bullet)$. By
Proposition~\ref{prop:refinement-equivariant}(4), the parent fibre
$(\pi_{n+1, n}^E)^{-1}(e)$ has exactly two oriented edges (the two
subdividing pieces of $e$), so the sum has two terms and the average
factor $\tfrac{1}{2}$ gives the standard arithmetic mean.

We verify that $p_{n+1, n}^1 A$ satisfies the antisymmetry
convention: $(p_{n+1, n}^1 A)(-e) = \tfrac{1}{2} \sum A(-e') =
-\tfrac{1}{2} \sum A(e') = -(p_{n+1, n}^1 A)(e)$, since each
subdividing piece of $-e$ is the reverse of a subdividing piece of
$e$.
\end{definition}

\subsection{Bonding maps are contractions}

\begin{theorem}[Bonding contractions]
\label{thm:bonding}
For every $n \geq 0$:
\begin{enumerate}
\item $p_{n+1, n}^0 \colon X_{n+1}^0 \to X_n^0$ is a contraction:
  $\|p_{n+1, n}^0 f\|_{X_n^0} \leq \|f\|_{X_{n+1}^0}$ for all
  $f \in X_{n+1}^0$.
\item $p_{n+1, n}^1 \colon X_{n+1}^1 \to X_n^1$ is a contraction:
  $\|p_{n+1, n}^1 A\|_{X_n^1} \leq \|A\|_{X_{n+1}^1}$ for all
  $A \in X_{n+1}^1$.
\end{enumerate}
\end{theorem}

\begin{proof}
\emph{(1)} By Definition~\ref{def:p0},
$\|p_{n+1, n}^0 f\|_{X_n^0}^2 = \sum_{v \in V(G_n^\bullet)}
\|f(v)\|_{F^0}^2$. Since $V(G_n^\bullet) \subset
V(G_{n+1}^\bullet)$ and $\|f(v)\|_{F^0}^2 \geq 0$ for every $v$,
\[
  \|p_{n+1, n}^0 f\|_{X_n^0}^2 = \sum_{v \in V(G_n^\bullet)}
    \|f(v)\|_{F^0}^2 \leq \sum_{w \in V(G_{n+1}^\bullet)}
    \|f(w)\|_{F^0}^2 = \|f\|_{X_{n+1}^0}^2.
\]

\emph{(2)} By Definition~\ref{def:p1} and the Cauchy--Schwarz
inequality applied to the two-term average $\tfrac{1}{2}(A(e_1) +
A(e_2))$:
\[
  \tfrac{1}{2}\|(p_{n+1, n}^1 A)(e)\|_{F^1}^2
  = \tfrac{1}{2} \cdot \tfrac{1}{4} \|A(e_1) + A(e_2)\|_{F^1}^2
  \leq \tfrac{1}{2} \cdot \tfrac{1}{2} (\|A(e_1)\|^2 + \|A(e_2)\|^2),
\]
where the first equality is the antisymmetry convention's
$\tfrac{1}{2}$ factor and the inequality is Cauchy--Schwarz on the
two-term average. Summing over $e \in E^{\rm or}(G_n^\bullet)$
and using that each subdividing piece $e'$ of level $n+1$
contributes to exactly one parent (by
Proposition~\ref{prop:refinement-equivariant}(4)):
\[
  \|p_{n+1, n}^1 A\|_{X_n^1}^2
  \leq \tfrac{1}{2} \sum_{e' \in E^{\rm or}_{\rm sub}(G_{n+1}^
    \bullet)} \|A(e')\|^2_{F^1}
  \leq \tfrac{1}{2} \sum_{e' \in E^{\rm or}(G_{n+1}^\bullet)}
    \|A(e')\|^2_{F^1}
  = \|A\|_{X_{n+1}^1}^2,
\]
where the final inequality is because $E^{\rm or}_{\rm sub}
(G_{n+1}^\bullet) \subseteq E^{\rm or}(G_{n+1}^\bullet)$.
\end{proof}

\subsection{Adjoint cylindrical embeddings}

\begin{definition}[Cylindrical embeddings]
\label{def:j}
For $n \geq 0$, define:
\begin{itemize}
\item \emph{Vertex embedding} $j_{n, n+1}^0 \colon X_n^0 \to
  X_{n+1}^0$ by zero-extension:
  \[
    (j_{n, n+1}^0 g)(w) := \begin{cases} g(w) & \text{if } w \in
      V(G_n^\bullet), \\ 0 & \text{if } w \in V(G_{n+1}^\bullet)
      \setminus V(G_n^\bullet). \end{cases}
  \]
\item \emph{Edge embedding (adjoint of $p^1$)}
  $j_{n, n+1}^1 \colon X_n^1 \to X_{n+1}^1$ by rescaled
  parent-replication:
  \[
    (j_{n, n+1}^1 B)(e') := \begin{cases}
      \tfrac{1}{2} B(\pi_{n+1, n}^E(e'))
        & \text{if } \pi_{n+1, n}^E(e') \neq \bot, \\
      0 & \text{otherwise.}
    \end{cases}
  \]
  The rescaling factor $\tfrac{1}{2}$ is required for $j_{n, n+1}^1$
  to be the genuine $\ell^2$-adjoint of $p_{n+1, n}^1$
  (Proposition~\ref{prop:adjoints}(3) below). Without this factor,
  the unscaled replication $s_{n, n+1}^1 B(e') := B(\pi^E(e'))$ is a
  \emph{section} of $p_{n+1, n}^1$ (i.e.\
  $p_{n+1, n}^1 \circ s_{n, n+1}^1 = \id$,
  Remark~\ref{rmk:edge-section} below) but not its adjoint.
\end{itemize}
\end{definition}

\begin{proposition}[Adjoint relation]
\label{prop:adjoints}
For every $n \geq 0$:
\begin{enumerate}
\item $j_{n, n+1}^0 = (p_{n+1, n}^0)^*$ as bounded linear maps,
  i.e.\ $\langle p_{n+1, n}^0 f, g \rangle_{X_n^0} = \langle f,
  j_{n, n+1}^0 g \rangle_{X_{n+1}^0}$ for all $f \in X_{n+1}^0,
  g \in X_n^0$.
\item $p_{n+1, n}^0 \circ j_{n, n+1}^0 = \id_{X_n^0}$ (vertex
  embedding is a section).
\item $j_{n, n+1}^1 = (p_{n+1, n}^1)^*$ as bounded linear maps.
\item $p_{n+1, n}^1 \circ j_{n, n+1}^1 = \tfrac{1}{2}\,\id_{X_n^1}$
  (the genuine adjoint $j^1$ is a half-section, not an exact
  section, of $p^1$; see also
  Remark~\ref{rmk:edge-section}).
\end{enumerate}
\end{proposition}

\begin{proof}
\emph{(1)}
\begin{align*}
  \langle p_{n+1, n}^0 f, g \rangle_{X_n^0}
  &= \sum_{v \in V(G_n^\bullet)} \langle f(v), g(v) \rangle_{F^0} \\
  &= \sum_{w \in V(G_{n+1}^\bullet)} \langle f(w), (j_{n, n+1}^0 g)
    (w) \rangle_{F^0} \quad \text{(extending by zero)} \\
  &= \langle f, j_{n, n+1}^0 g \rangle_{X_{n+1}^0}.
\end{align*}

\emph{(2)} For $g \in X_n^0$ and $v \in V(G_n^\bullet)$,
$(p_{n+1, n}^0 j_{n, n+1}^0 g)(v) = (j_{n, n+1}^0 g)(v) = g(v)$.

\emph{(3)} Compute both sides explicitly. With the edge-sector
inner product factor of $\tfrac{1}{2}$ from
Definition~\ref{def:Xn-1}:
\begin{align*}
  \langle p_{n+1, n}^1 A, B \rangle_{X_n^1}
  &= \tfrac{1}{2} \sum_{e \in E^{\rm or}(G_n^\bullet)}
    \langle (p_{n+1, n}^1 A)(e), B(e) \rangle_{F^1} \\
  &= \tfrac{1}{2} \sum_{e} \tfrac{1}{2} \sum_{e' \in (\pi^E)^{-1}(e)}
    \langle A(e'), B(e) \rangle_{F^1} \\
  &= \tfrac{1}{4} \sum_{e' \in E^{\rm or}_{\rm sub}(G_{n+1}^
    \bullet)} \langle A(e'), B(\pi^E(e')) \rangle_{F^1}.
\end{align*}
With the rescaled $j_{n, n+1}^1$ formula
$(j_{n, n+1}^1 B)(e') = \tfrac{1}{2} B(\pi^E(e'))$ on subdividing
edges and $0$ elsewhere:
\begin{align*}
  \langle A, j_{n, n+1}^1 B \rangle_{X_{n+1}^1}
  &= \tfrac{1}{2} \sum_{e' \in E^{\rm or}(G_{n+1}^\bullet)}
    \langle A(e'), (j_{n, n+1}^1 B)(e') \rangle_{F^1} \\
  &= \tfrac{1}{2} \sum_{e' \in E^{\rm or}_{\rm sub}}
    \langle A(e'), \tfrac{1}{2} B(\pi^E(e')) \rangle_{F^1} \\
  &= \tfrac{1}{4} \sum_{e' \in E^{\rm or}_{\rm sub}}
    \langle A(e'), B(\pi^E(e')) \rangle_{F^1}.
\end{align*}
The two expressions agree, so $j_{n, n+1}^1 = (p_{n+1, n}^1)^*$.

\emph{(4)} For $B \in X_n^1$ and $e \in E^{\rm or}(G_n^\bullet)$,
the parent fibre of $e$ has two oriented edges $e_1, e_2$, both
satisfying $\pi^E(e_i) = e$. Thus
$(j_{n, n+1}^1 B)(e_i) = \tfrac{1}{2} B(e)$ for $i = 1, 2$, and
$(p_{n+1, n}^1 j_{n, n+1}^1 B)(e) = \tfrac{1}{2}\bigl((j_{n, n+1}^1
B)(e_1) + (j_{n, n+1}^1 B)(e_2)\bigr) = \tfrac{1}{2}\bigl(\tfrac{1}{2}
B(e) + \tfrac{1}{2} B(e)\bigr) = \tfrac{1}{2} B(e)$. So
$p_{n+1, n}^1 \circ j_{n, n+1}^1 = \tfrac{1}{2}\,\id_{X_n^1}$.
\end{proof}

\begin{remark}[Edge section vs adjoint]
\label{rmk:edge-section}
The edge sector has \emph{two} naturally-related upward maps from
$X_n^1$ to $X_{n+1}^1$:
\begin{itemize}
\item the \emph{adjoint} $j_{n, n+1}^1 := (p_{n+1, n}^1)^*$, which
  satisfies $j_{n, n+1}^1 B(e') = \tfrac{1}{2} B(\pi^E(e'))$ on
  subdividing edges and $p_{n+1, n}^1 \circ j_{n, n+1}^1 =
  \tfrac{1}{2} \id$ (Proposition~\ref{prop:adjoints}(4));
\item the \emph{section} $s_{n, n+1}^1 B(e') := B(\pi^E(e'))$
  on subdividing edges and $0$ elsewhere, which satisfies
  $p_{n+1, n}^1 \circ s_{n, n+1}^1 = \id$ but is not the
  $\ell^2$-adjoint of $p_{n+1, n}^1$.
\end{itemize}
The two differ by the rescaling $j_{n, n+1}^1 = \tfrac{1}{2}
s_{n, n+1}^1$. We use $j$ (the adjoint) below for inverse-limit
constructions; we use $s$ (the section) only when an exact
section is needed downstream. The vertex sector does not have this
distinction because the vertex bonding map is restriction (not
averaging), and the corresponding adjoint $j_{n, n+1}^0$ is
zero-extension, which is a section ($p^0 \circ j^0 = \id$).
\end{remark}

%=====================================================================
\section{The Cylindrical Algebraic Direct-Limit Space}
\label{sec:cyl-lim}
%=====================================================================

\subsection{Algebraic direct limit}

\begin{definition}[Algebraic direct limit on vertices]
\label{def:Xcyl-vertex}
The \emph{cylindrical algebraic direct-limit space (vertex sector)}
is
\[
  X_\cyl^{0, \bullet} := \varinjlim_n \big( X_n^{0, \bullet},\,
    j_{n, n+1}^0 \big),
\]
the direct limit of the system $(X_n^{0, \bullet})_{n \geq 0}$
under the cylindrical embeddings $j_{n, n+1}^0$ of
Definition~\ref{def:j}.
\end{definition}

\begin{theorem}[Identification of $X_\cyl^0$]
\label{thm:Xcyl-direct-limit}
The vertex cylindrical direct-limit space identifies canonically
with $c_{00}$-functions on the countable union $V_\infty^\bullet$
of refinement-vertex sets:
\[
  X_\cyl^{0, \bullet} \cong c_{00}\bigl(V_\infty^\bullet;\, F^0\bigr),
\]
where $V_\infty^\bullet := \bigcup_{n \geq 0} V(G_n^\bullet)$
(a countable set, by Proposition~\ref{prop:refinement-equivariant}(1)
and the inclusion chain) and $c_{00}(\cdot; F^0)$ denotes
$F^0$-valued functions of finite support.
\end{theorem}

\begin{proof}
The algebraic direct limit $\varinjlim_n X_n^{0, \bullet}$ is the
disjoint union $\bigsqcup_n X_n^{0, \bullet}$ modulo the
equivalence relation generated by $f \sim j_{n, n+1}^0 f$ for
$f \in X_n^{0, \bullet}$ \cite[Ch.~II.1.4]{ReedSimonI}. Each
equivalence class has a unique representative at the smallest
level $n$ for which the class contains an element; that
representative is a function on $V(G_n^\bullet)$, which by
zero-extension to all higher levels gives a function on
$V_\infty^\bullet$ vanishing outside $V(G_n^\bullet)$.

Conversely, every $c_{00}$-function $\phi \colon V_\infty^\bullet
\to F^0$ has support contained in $V(G_n^\bullet)$ for some
$n$ (since the support is finite and $V_\infty^\bullet$ is the
ascending union), so it represents an element of $X_n^{0, \bullet}$.

These maps are inverse to each other and are $\R$-linear, giving
the identification $X_\cyl^{0, \bullet} \cong c_{00}(V_\infty^
\bullet; F^0)$.
\end{proof}

\begin{remark}[Edge sector: half-section, finite-level only]
\label{rmk:edge-cyl}
The edge cylindrical direct-limit construction of the vertex
sector does \emph{not} transfer cleanly to edges. The adjoint
$j_{n, n+1}^1 = (p_{n+1, n}^1)^*$ satisfies only
$p_{n+1, n}^1 \circ j_{n, n+1}^1 = \tfrac{1}{2}\id_{X_n^1}$
(Proposition~\ref{prop:adjoints}(4)): it is a \emph{half-section},
not a section, of $p_{n+1, n}^1$. Consequently
$\varinjlim_n(X_n^{1, \bullet}, j_{n, n+1}^1)$ does not enjoy the
properties used in the vertex identification
(Theorem~\ref{thm:Xcyl-direct-limit}), and we do \emph{not} define an edge
cylindrical direct-limit space here. Downstream papers requiring
the edge sector should work at the finite level $X_n^1$ with the
bounded bonding maps $p_{n+1, n}^1$ and their adjoints, without
relying on a cylindrical-direct-limit formalism for edges.
\end{remark}

\subsection{The $\ell^2$-norm on $X_\cyl$}

\begin{definition}[Pre-Hilbert structure on $X_\cyl$]
\label{def:Xcyl-norm}
On $X_\cyl^{0, \bullet} \cong c_{00}(V_\infty^\bullet; F^0)$, define
the inner product
\[
  \langle \phi, \psi \rangle_{X_\cyl^0} :=
    \sum_{v \in V_\infty^\bullet} \langle \phi(v), \psi(v)
    \rangle_{F^0}.
\]
This sum is finite since $\phi, \psi$ have finite support. The
resulting inner product is positive-definite (since $\langle \cdot,
\cdot \rangle_{F^0}$ is) and makes $X_\cyl^0$ into a pre-Hilbert
space.
\end{definition}

\begin{lemma}[Compatibility with finite-level inner products]
\label{lem:Xcyl-compatibility}
For each $n \geq 0$ and $f \in X_n^{0, \bullet}$ identified with
its image in $X_\cyl^{0, \bullet}$ via the embedding $j_{n,
\infty}^0$ (zero-extension to $V_\infty^\bullet$),
$\|f\|_{X_n^0} = \|f\|_{X_\cyl^0}$.
\end{lemma}

\begin{proof}
$\|f\|_{X_n^0}^2 = \sum_{v \in V(G_n^\bullet)} \|f(v)\|^2_{F^0}$;
$\|f\|_{X_\cyl^0}^2 = \sum_{v \in V_\infty^\bullet} \|f(v)\|^2$,
where $f(v) = 0$ for $v \in V_\infty^\bullet \setminus V(G_n^
\bullet)$. The two sums agree.
\end{proof}

%=====================================================================
\section{The Inverse-Limit Hilbert Space}
\label{sec:inv-lim}
%=====================================================================

\begin{definition}[Inverse-limit space]
\label{def:Xinv}
The \emph{vertex inverse-limit space} is
\[
  X_\inv^{0, \bullet} :=
    \big\{(f_n)_{n \geq 0} : f_n \in X_n^{0, \bullet},\
    p_{n+1, n}^0 f_{n+1} = f_n \text{ for all } n,\
    \sup_n \|f_n\|_{X_n^0} < \infty \big\}.
\]
It is equipped with the inner product
\[
  \langle f, g \rangle_{X_\inv^0} := \lim_{n \to \infty}
    \langle f_n, g_n \rangle_{X_n^0}.
\]
\end{definition}

\begin{lemma}[Limits exist]
\label{lem:Xinv-limit-exists}
For every bounded compatible family $(f_n)_n \in X_\inv^{0, \bullet}$,
the sequence $\|f_n\|_{X_n^0}$ is non-decreasing and bounded, hence
converges. By polarization, $\langle f_n, g_n \rangle_{X_n^0}$
also converges for any pair of compatible families $(f_n), (g_n)$.
\end{lemma}

\begin{proof}
By Definition~\ref{def:p0}, $p_{n+1, n}^0$ is restriction to
$V(G_n^\bullet) \subset V(G_{n+1}^\bullet)$. So the compatibility
$f_n = p_{n+1, n}^0 f_{n+1}$ means $f_n$ is the restriction of
$f_{n+1}$ to the smaller domain $V(G_n^\bullet)$. Hence
$\|f_n\|^2 = \sum_{v \in V(G_n^\bullet)} \|f_{n+1}(v)\|^2 \leq
\sum_{w \in V(G_{n+1}^\bullet)} \|f_{n+1}(w)\|^2 = \|f_{n+1}\|^2$.
The sequence is non-decreasing. By the supremum bound it is
bounded, hence converges. Polarization for cross terms:
$\langle f_n, g_n \rangle = \tfrac{1}{4}(\|f_n + g_n\|^2 -
\|f_n - g_n\|^2)$, both norms converge, so does the inner product.
\end{proof}

\begin{theorem}[$X_\inv^0 \cong \ell^2(V_\infty^\bullet)$]
\label{thm:Xinv-Hilbert}
The map
\[
  \Phi \colon X_\inv^{0, \bullet} \to \ell^2(V_\infty^\bullet;
    F^0), \qquad
  \Phi\bigl((f_n)_n\bigr) := f_\infty,
\]
where $f_\infty(v) := f_n(v)$ for $v \in V(G_n^\bullet)$ (using
the compatibility $f_n = p_{n+1, n}^0 f_{n+1}$ to ensure
well-definedness across levels), is an $\R$-linear bijection that
intertwines the inner products:
\[
  \langle f, g \rangle_{X_\inv^0} = \langle \Phi(f), \Phi(g)
    \rangle_{\ell^2(V_\infty^\bullet; F^0)}.
\]
In particular, $X_\inv^{0, \bullet}$ is a separable Hilbert space.
\end{theorem}

\begin{proof}
\emph{Well-definedness of $\Phi$.} For $v \in V_\infty^\bullet =
\bigcup_n V(G_n^\bullet)$, choose $n$ minimal such that $v \in
V(G_n^\bullet)$ and set $f_\infty(v) := f_n(v)$. The compatibility
$f_n = p_{n+1, n}^0 f_{n+1}$ gives $f_n(v) = f_{n+1}(v)$ for
$v \in V(G_n^\bullet)$ at any larger level, so the value is
independent of the choice of $n$.

$f_\infty \in \ell^2(V_\infty^\bullet; F^0)$: by
Lemma~\ref{lem:Xinv-limit-exists}, $\sum_v \|f_\infty(v)\|^2 =
\lim_n \sum_{v \in V(G_n^\bullet)} \|f_n(v)\|^2 = \lim_n
\|f_n\|^2 < \infty$.

\emph{Injectivity.} If $\Phi(f) = 0$, then $f_\infty(v) = 0$ for
all $v$, so $f_n(v) = 0$ for all $n, v \in V(G_n^\bullet)$, hence
$f_n = 0$ for all $n$.

\emph{Surjectivity.} For $\phi \in \ell^2(V_\infty^\bullet; F^0)$,
define $f_n := \phi|_{V(G_n^\bullet)} \in X_n^0$. Then
$\|f_n\|^2 = \sum_{v \in V(G_n^\bullet)} \|\phi(v)\|^2 \leq
\|\phi\|^2_{\ell^2}$, bounded. Compatibility: $p_{n+1, n}^0
f_{n+1} = (\phi|_{V(G_{n+1}^\bullet)})|_{V(G_n^\bullet)} =
\phi|_{V(G_n^\bullet)} = f_n$. So $(f_n) \in X_\inv$ and
$\Phi(f) = \phi$.

\emph{Inner product preservation.}
$\langle f, g \rangle_{X_\inv} = \lim_n \sum_{v \in V(G_n^
\bullet)} \langle f_n(v), g_n(v) \rangle = \sum_{v \in V_\infty^
\bullet} \langle \Phi(f)(v), \Phi(g)(v) \rangle =
\langle \Phi(f), \Phi(g) \rangle_{\ell^2}$.

\emph{Hilbert + separable.} $\ell^2$ on a countable set with
finite-dimensional fibre is a standard separable Hilbert space
\cite[Ch.~II.4]{ReedSimonI}. The bijection $\Phi$ transports
this structure to $X_\inv^{0, \bullet}$.
\end{proof}

\subsection{Density of $X_\cyl$ in $X_\inv$}

\begin{proposition}[$X_\cyl^0$ is dense in $X_\inv^0$]
\label{prop:Xcyl-dense}
Under the canonical map
\[
  X_\cyl^{0, \bullet} \cong c_{00}(V_\infty^\bullet; F^0)
    \hookrightarrow \ell^2(V_\infty^\bullet; F^0) \cong
    X_\inv^{0, \bullet},
\]
the cylindrical algebraic direct-limit space embeds as a dense
subspace of the inverse-limit Hilbert space.
\end{proposition}

\begin{proof}
$c_{00}$ is dense in $\ell^2$ on any countable index set, by
truncation: for $\phi \in \ell^2(V_\infty^\bullet)$ and
$\varepsilon > 0$, choose $n$ large enough that $\sum_{v \in
V_\infty^\bullet \setminus V(G_n^\bullet)} \|\phi(v)\|^2 <
\varepsilon^2$, then $\phi_n := \phi|_{V(G_n^\bullet)}$ extended
by zero is in $c_{00}$ and satisfies $\|\phi - \phi_n\|^2 <
\varepsilon^2$. By Theorem~\ref{thm:Xinv-Hilbert}, the same
density holds in $X_\inv^{0, \bullet}$.
\end{proof}

\subsection{Edge sector}

\begin{remark}[Edge inverse limit: deferred]
\label{rmk:edge-inv}
The edge bonding map $p_{n+1, n}^1$ is parent-fibre averaging
(not restriction), and satisfies the contraction bound
$\|p_{n+1,n}^1\| \leq 1$ (Theorem~\ref{thm:bonding}(2)). However,
contraction-monotonicity alone does \emph{not} establish completeness
of the set
$X_\inv^{1, \bullet} := \{(A_n) : p_{n+1, n}^1 A_{n+1} = A_n,\
\sup_n \|A_n\| < \infty\}$
in the induced limit norm: the vertex proof of
Theorem~\ref{thm:Xinv-Hilbert} uses the concrete identification
with $\ell^2(V_\infty^\bullet; F^0)$ (via restriction-bonding maps),
and the averaging-bonding edge sector has no analogous
identification at present. We therefore do \emph{not} claim here
that $X_\inv^{1, \bullet}$ is a Hilbert space; it is at this stage
a set of bounded compatible families with a pre-inner product
induced by the finite-level norms. Establishing Hilbert-completion
for the edge sector is left to downstream papers needing it.
\end{remark}

%=====================================================================
\section{Commuting Symmetry Actions}
\label{sec:commute}
%=====================================================================

\subsection{The three actions on $X_n$}

\begin{definition}[Coxeter action, branchwise]
\label{def:coxeter}
For $g \in W^\bullet$ and $f \in X_n^{0, \bullet}$, define
$(g \cdot f)(v) := \rho_{F^0}^\bullet(g) \cdot f(g^{-1} \cdot v)$
for $v \in V(G_n^\bullet)$, where the fibre representation
$\rho_{F^0}^\bullet$ depends on the branch:
\begin{itemize}
\item On the $H_4$ branch ($\bullet = H_4$, $W^\bullet = W(H_4)$):
  $\rho_{F^0}^{H_4}(g)$ acts by the standard $H_4$-reflection
  representation on $V_{H_4}$, and trivially on $V_{D_4},
  V_{16}, V_8$.
\item On the $D_4$ branch ($\bullet = D_4$, $W^\bullet = W(D_4)$):
  $\rho_{F^0}^{D_4}(g)$ acts by the standard $D_4$-reflection
  representation on $V_{D_4}$, and trivially on $V_{H_4},
  V_{16}, V_8$.
\end{itemize}
The analogous action on $X_n^{1, \bullet}$ acts trivially on $F^1
= \Im(\OO)$ on both branches and on the edge index by $W^\bullet$-
permutation.
\end{definition}

\begin{lemma}[Unitarity of the Coxeter action]
\label{lem:coxeter-unitary}
Under the $\ell^2$-inner product of Definition~\ref{def:Xn-0}, the
Coxeter action of Definition~\ref{def:coxeter} is unitary on
$X_n^{0, \bullet}$ and on $X_n^{1, \bullet}$.
\end{lemma}

\begin{proof}
The $H_4$- (resp.\ $D_4$-) reflection representation on $V_{H_4}$
(resp.\ $V_{D_4}$) is orthogonal with respect to the Euclidean
metric on the fibre (P2~\S\S 3--5), and acts trivially on the other
fibre summands; hence $\rho_{F^0}^\bullet(g)$ is a fibre isometry.
The vertex (resp.\ oriented-edge) measure $\mu_n^\bullet$ (resp.\
$\nu_n^\bullet$) is $W^\bullet$-invariant by the
Coxeter-equivariance of the refinement (Proposition~\ref
{prop:refinement-equivariant}). The combined pre-composition with
$g^{-1}$ and fibrewise action by $\rho_{F^0}^\bullet(g)$ is
therefore a unitary on the $\ell^2$-space.
\end{proof}

\subsection{$\Q$-form and the $\sigma$-action}

\begin{definition}[Mixed $\Q(\varphi)/\Q$-form of the vertex sector]
\label{def:Xn-Q-form}
By inspection of the explicit fibre constructions in
\cite[\S\S 3, 5, 6, 7]{AlgebraicSubstrate}: $V_{D_4} = \R^4$ has
the standard $\Q$-form $V_{D_4, \Q} := \Q^4$; $V_{16} =
\Cl(1, 3)$ has the $\Q$-form $V_{16, \Q} := \Q\{\gamma_A\}$
(rational span of the Clifford monomial basis,
\cite[Definition~\texttt{def:Cl-1-3-basis}]{AlgebraicSubstrate}, which is a $\Q$-basis of a
$\Q$-form whose $\R$-extension equals $\Cl(1, 3)$); $V_{8} = \OO$
has the $\Q$-form $V_{8, \Q} := \Q\{1, e_1, \ldots, e_7\}$
(rational span of the standard octonion basis,
\cite[Definition~\texttt{def:octonion-basis}]{AlgebraicSubstrate}). For $V_{H_4} = \R^4$ the
natural rational structure is the $\Q(\varphi)$-form
$V_{H_4, \Q(\varphi)} := \Q(\varphi)^4 \subset \R^4$ via the
standard real embedding $\iota \colon \Q(\varphi) \hookrightarrow \R$
sending $\varphi$ to the golden ratio
$\tfrac12(1 + \sqrt{5})$, because the icosian arithmetic on
$V_{H_4}$ is naturally over $\Z[\varphi]$ rather than over $\Z$
\cite[\S 3]{AlgebraicSubstrate}.

The combined rationality data is the \emph{mixed
$\Q(\varphi)/\Q$-form}
\[
  F^0_{\Q(\varphi)/\Q} := V_{H_4, \Q(\varphi)} \oplus V_{D_4, \Q}
    \oplus V_{16, \Q} \oplus V_{8, \Q},
\]
which is a $\Q$-vector space of total $\Q$-dimension
$8 + 4 + 16 + 8 = 36$ (since $V_{H_4, \Q(\varphi)} = \Q(\varphi)^4$
has $\Q$-dimension $8$). Its R-extension is sector-dependent: on the
$H_4$ sector we use $V_{H_4, \Q(\varphi)} \otimes_{\Q(\varphi)} \R
= \R^4 = V_{H_4}$ via $\iota$, on the other three sectors we use
$V_{*, \Q} \otimes_\Q \R = V_*$. With this convention
\[
  F^0 \;=\; V_{H_4} \oplus V_{D_4} \oplus V_{16} \oplus V_{8}
  \quad (\dim_\R F^0 = 32)
\]
is recovered from $F^0_{\Q(\varphi)/\Q}$. We do \emph{not} claim
$F^0_{\Q(\varphi)/\Q} \otimes_\Q \R = F^0$; that would have
$\R$-dimension $36$, not $32$, because $\Q(\varphi) \otimes_\Q \R
\cong \R \oplus \R$ (the two real roots of $x^2 - x - 1$). The
$\R$-dimension $32$ of $F^0$ comes from the $\Q(\varphi)$-tensor
product on the $H_4$ sector, not from the $\Q$-tensor product.

Define the mixed-form vertex space at level $n$ by
\[
  X_n^{0, \bullet}|_{\Q(\varphi)/\Q}
  := \ell^2_{\rm fin}\bigl(V(G_n^\bullet),
    \mu_n^\bullet;\, F^0_{\Q(\varphi)/\Q}\bigr),
\]
the $F^0_{\Q(\varphi)/\Q}$-valued functions on the finite set
$V(G_n^\bullet)$.
\end{definition}

\begin{definition}[$\sigma$-action on the mixed-form vertex sector]
\label{def:sigma}
The Galois involution $\sigma \colon \Q(\varphi) \to \Q(\varphi)$
defined by $\sigma(\varphi) = 1 - \varphi$
(\cite[Def.~\texttt{def:sigma}]{SigmaRationality}) extends
sector-by-sector to a $\Q$-linear involution
$\sigma_n \colon X_n^{0, \bullet}|_{\Q(\varphi)/\Q}
\to X_n^{0, \bullet}|_{\Q(\varphi)/\Q}$
on the mixed-form vertex space of
Definition~\ref{def:Xn-Q-form} as follows:
\begin{itemize}[topsep=2pt]
\item on the $V_{H_4, \Q(\varphi)} = \Q(\varphi)^4$ component,
  $\sigma_n$ acts coefficientwise by $\sigma$ in each of the four
  $\Q(\varphi)$-coordinates;
\item on each of $V_{D_4, \Q}, V_{16, \Q}, V_{8, \Q}$, $\sigma_n$
  acts as the identity (these sectors carry no
  $\Q(\varphi)$-rationality).
\end{itemize}
$\sigma_n$ is a $\Q$-linear involution on the $\Q$-vector space
$X_n^{0, \bullet}|_{\Q(\varphi)/\Q}$. It does \emph{not} extend to
an $\R$-linear involution on the $\R$-extension
$X_n^{0, \bullet} = \ell^2(V(G_n^\bullet), \mu_n^\bullet; F^0)$
through the standard real embedding $\iota$ (because $\sigma$ on
$\Q(\varphi) \otimes_{\Q(\varphi), \iota} \R = \R$ is fixed by
$\iota$ but free over the conjugate embedding $\bar\iota$ sending
$\varphi$ to $1 - \varphi$); $\sigma_n$ is the \emph{rationality-side}
involution, used in the intertwining identities below at the
mixed-form level.

We do not define a $\sigma$-action on the edge sector $X_n^1$ in
this paper: the edge fibre $\Im(\OO)$ has no canonical
$\Q(\varphi)$-coefficient structure relevant to the cascade
arithmetic, and downstream papers needing edge $\sigma$-actions
construct them separately.
\end{definition}

\subsection{The intertwining theorem}

\begin{theorem}[Intertwining identities]
\label{thm:commute}
For every $n \geq 0$ and $\bullet \in \{H_4, D_4\}$:
\begin{enumerate}
\item (Refinement-Coxeter, vertex sector.) For every $g \in W^
  \bullet$,
  \[
    g \circ p_{n+1, n}^0 = p_{n+1, n}^0 \circ g
    \quad \text{on } X_{n+1}^{0, \bullet},
  \]
  and analogously $g \circ j_{n, n+1}^0 = j_{n, n+1}^0 \circ g$.
\item (Refinement-Coxeter, edge sector.) The same identities hold
  for $p_{n+1, n}^1$ and $j_{n, n+1}^1$ on $X_{n+1}^{1, \bullet}$.
\item (Refinement-$\sigma$, vertex sector, mixed-form level.) On
  the mixed-form vertex space
  $X_{n+1}^{0, \bullet}|_{\Q(\varphi)/\Q}$ of
  Definition~\ref{def:Xn-Q-form},
  \[
    \sigma_{n} \circ p_{n+1, n}^0|_{\Q(\varphi)/\Q}
    = p_{n+1, n}^0|_{\Q(\varphi)/\Q} \circ \sigma_{n+1},
  \]
  where $p_{n+1, n}^0|_{\Q(\varphi)/\Q}$ is the mixed-form
  restriction of the bonding map (well-defined because vertex
  restriction preserves the sectorwise rationality)
  and $\sigma_n, \sigma_{n+1}$ are the involutions of
  Definition~\ref{def:sigma}. Analogously, with $j_{n, n+1}^0$
  the (zero-extension) adjoint of $p_{n+1, n}^0$ from
  Proposition~\ref{prop:adjoints},
  \[
    \sigma_{n+1} \circ j_{n, n+1}^0|_{\Q(\varphi)/\Q}
    = j_{n, n+1}^0|_{\Q(\varphi)/\Q} \circ \sigma_n
  \]
  on $X_n^{0, \bullet}|_{\Q(\varphi)/\Q}$ (the $\sigma$ indices swap
  relative to the $p^0$ identity since $j$ goes upward
  $n \to n+1$ while $p$ goes downward).
\item (Sectorwise Coxeter-$\sigma$, vertex sector.) On the $D_4$
  branch ($\bullet = D_4$, $W^\bullet = W(D_4)$), the action of
  Definition~\ref{def:coxeter} acts on $V_{D_4}$ and trivially on
  $V_{H_4}, V_{16}, V_8$; since $\sigma$ acts trivially on
  $V_{D_4, \Q}, V_{16, \Q}, V_{8, \Q}$, the commutation
  $g|_{\Q(\varphi)} \circ \sigma_n = \sigma_n \circ g|_{\Q(\varphi)}$
  holds automatically for all $g \in W(D_4)$ on
  $X_n^{0, D_4}|_{\Q(\varphi)}$. On the $H_4$ branch, the analogous
  automatic-commutation statement holds only on the
  $\sigma$-trivial summands $V_{D_4, \Q}, V_{16, \Q}, V_{8, \Q}$
  (on which $W(H_4)$ acts trivially by Definition~\ref{def:coxeter}).
  The full statement on the $V_{H_4}$ summand is \emph{not} an
  instance of \cite[Theorem~\texttt{thm:sigma-commutes}]{SigmaRationality}
  (the source theorem requires endomorphisms defined over the
  $\sigma$-fixed $\Q$-form, which $W(H_4)$ reflections with
  $\Q(\varphi)$-rational coefficients are not), and is not claimed
  here; see Remark~\ref{rmk:coxeter-sigma-restricted}.
\end{enumerate}
The Coxeter intertwinings of (1)--(2) descend to the vertex
cylindrical and inverse-limit spaces $X_\cyl^{0, \bullet}$,
$X_\inv^{0, \bullet}$ by linearity and continuity (using
Theorem~\ref{thm:Xinv-Hilbert} for the inverse limit). The
$\sigma$-intertwining of (3) is stated only at the mixed-form level
$X_n^{0, \bullet}|_{\Q(\varphi)/\Q}$; we do not claim a
$\sigma$-action on the $\R$-Hilbert spaces $X_n^{0, \bullet}$,
$X_\cyl^{0, \bullet}$, $X_\inv^{0, \bullet}$, since (per
Definition~\ref{def:sigma}) $\sigma$ does not extend $\R$-linearly
through the standard real embedding $\iota$.
\end{theorem}

\begin{remark}[Why Coxeter-$\sigma$ commutation is restricted]
\label{rmk:coxeter-sigma-restricted}
\cite[Theorem~\texttt{thm:sigma-commutes}]{SigmaRationality}
requires a $\Q$-linear endomorphism $T$ of the $\Q$-vector space
$V$ \emph{defined over the $\sigma$-fixed $\Q$-form}: concretely,
$T$ must restrict to $V \otimes 1 \subset V \otimes_\Q \Q(\varphi)$.
On the $V_{H_4, \Q(\varphi)}$ component, the $W(H_4)$ action is
given by reflection matrices whose coefficients lie in
$\Q(\varphi)$ (not in $\Q$), because the icosian roots have
$\Q(\varphi)$-rational coordinates
\cite{AlgebraicSubstrate}. Such a reflection is
$\Q(\varphi)$-linear (hence a fortiori $\Q$-linear on the
underlying $\Q$-vector space $V_{H_4, \Q(\varphi)}$), but it does
\emph{not} restrict to a $\Q$-linear endomorphism of any $\Q$-form
$V_{H_4, \Q} \subset V_{H_4, \Q(\varphi)}$ (e.g.\ $\Q^4$ via the
standard basis): a non-trivial $\Q(\varphi)$-coefficient reflection
of a $\Q^4$-vector carries it out of $\Q^4$ into the general
$\Q(\varphi)^4$ space. The hypothesis of
\cite[Theorem~\texttt{thm:sigma-commutes}]{SigmaRationality} is
therefore not satisfied in the form required for automatic
$\sigma$-commutation, so the source theorem does not deliver
$g \circ \sigma_n = \sigma_n \circ g$ on the $V_{H_4}$ sector. We
do not assert in this paper that the equality fails; we assert
only that it is not delivered automatically by the source theorem
and is not claimed here.

The sectorwise statement (4) records the cases where commutation
\emph{does} hold automatically under the branchwise action of
Definition~\ref{def:coxeter}: the entire $D_4$-branch action (since
$W(D_4)$ acts only on $V_{D_4}$, which is $\sigma$-trivial), and
the $H_4$-branch action restricted to the $\sigma$-trivial summands
$V_{D_4, \Q}, V_{16, \Q}, V_{8, \Q}$ (on which $W(H_4)$ acts
trivially).

Downstream papers needing the full $W(H_4)$-$\sigma$ relationship
(e.g.\ in spectral-arithmetic constructions) should formulate it
as an intertwining with $\sigma$-conjugate Coxeter elements
$g^\sigma$ rather than as commutation: $\sigma \circ g = g^\sigma
\circ \sigma$ where $g^\sigma$ is the matrix-coefficient
$\sigma$-conjugate of $g$.
\end{remark}

\begin{proof}
\emph{(1) Refinement-Coxeter, vertex.} For $f \in X_{n+1}^0,
v \in V(G_n^\bullet), g \in W^\bullet$:
\begin{align*}
  (g \circ p_{n+1, n}^0 f)(v)
  &= \rho_{F^0}(g) \cdot (p_{n+1, n}^0 f)(g^{-1} \cdot v)
  = \rho_{F^0}(g) \cdot f(g^{-1} \cdot v) \\
  &= (g \cdot f)(v) = (p_{n+1, n}^0 (g \cdot f))(v).
\end{align*}
This uses $g^{-1} \cdot v \in V(G_n^\bullet)$ (because $V(G_n^
\bullet)$ is $W^\bullet$-stable) and the equality $f(g^{-1} v) =
(p_{n+1, n}^0 f)(g^{-1} v)$ (definition of $p^0$ as restriction).

The $j$-version follows by adjunction: $g \circ j = j \circ g$ is
equivalent to $g^{-1} \circ p = p \circ g^{-1}$ (which we just
proved with $g^{-1}$ in place of $g$).

\emph{(2) Refinement-Coxeter, edge.} For $A \in X_{n+1}^1,
e \in E^{\rm or}(G_n^\bullet), g \in W^\bullet$ (with the $W^\bullet$
action on $\Im(\OO)$ trivial, so we drop the $\rho_{F^1}(g)$ factor):
\begin{align*}
  (g \circ p_{n+1, n}^1 A)(e)
  &= (p_{n+1, n}^1 A)(g^{-1} \cdot e)
  = \tfrac{1}{2} \sum_{e' \in (\pi^E)^{-1}(g^{-1} e)} A(e').
\end{align*}
By Proposition~\ref{prop:refinement-equivariant}(3),
$(\pi^E)^{-1}(g^{-1} e) = g^{-1} \cdot (\pi^E)^{-1}(e)$.
Substituting $e'' := g \cdot e'$:
\[
  = \tfrac{1}{2} \sum_{e'' \in (\pi^E)^{-1}(e)}
    A(g^{-1} \cdot e'') = \tfrac{1}{2} \sum_{e''} (g \cdot A)(e'')
  = (p_{n+1, n}^1 (g \cdot A))(e).
\]

\emph{(3) Refinement-$\sigma$.} The bonding map $p_{n+1, n}^0$
restricts to a $\Q$-linear map $p_{n+1, n}^0|_\Q \colon
X_{n+1}^{0, \bullet}|_\Q \to X_n^{0, \bullet}|_\Q$ between
$\Q$-vector spaces of finite $\Q$-dimensions
$|V(G_{n+1}^\bullet)| \cdot 36$ and $|V(G_n^\bullet)| \cdot 36$
(using the $\Q$-dimension of $F^0_\Q$ from
Definition~\ref{def:Xn-Q-form}; the corresponding $\R$-dimensions
of the unrestricted spaces are
$|V(\cdot)| \cdot 32$, but the $\Q$-form has $\Q$-dimension $36$
because of the $\Q(\varphi)$-coefficient pair on the $V_{H_4}$
sector). By the general
functoriality theorem \cite[Theorem~\texttt{thm:sigma-commutes}]{SigmaRationality}, the scalar
extension $p_{n+1, n}^0|_\Q \otimes \id_{\Q(\varphi)}$ commutes
with the coefficientwise $\sigma$-actions on both sides.
The same argument applies to $j_{n, n+1}^0|_\Q$.

\emph{(4) Sectorwise Coxeter-$\sigma$.}
On the $D_4$ branch, $W(D_4)$ acts (by Definition~\ref{def:coxeter})
on $V_{D_4}$ and trivially on $V_{H_4}, V_{16}, V_8$; since $\sigma$
acts trivially on $V_{D_4, \Q}$ and on $V_{16, \Q}, V_{8, \Q}$, the
composition $g \circ \sigma$ equals $\sigma \circ g$ on every fibre
summand, so it equals $\sigma \circ g$ on
$X_n^{0, D_4}|_{\Q(\varphi)}$. On the $H_4$ branch, $W(H_4)$ acts
on $V_{H_4}$ and trivially on $V_{D_4}, V_{16}, V_8$; on the latter
three summands $\sigma$ also acts trivially, so the commutation is
automatic there by the same fibrewise computation. The remaining
$V_{H_4}$ summand is excluded from (4)
(Remark~\ref{rmk:coxeter-sigma-restricted}).

The descent to $X_\cyl$ and $X_\inv$: each action is continuous on
each $X_n$, and the intertwining is preserved under the algebraic
direct-limit construction (Theorem~\ref{thm:Xcyl-direct-limit})
and the $\ell^2$-completion (Theorem~\ref{thm:Xinv-Hilbert}).
\end{proof}

%=====================================================================
\section{Citation Contract and Downstream Handoff}
\label{sec:handoff}
%=====================================================================

\subsection{What downstream papers may cite}

\begin{enumerate}
\item Definitions~\ref{def:Xn-0}, \ref{def:Xn-1}: finite-level
  Hilbert sector spaces $X_n^{0, \bullet}, X_n^{1, \bullet}$
  under counting measure.
\item Theorem~\ref{thm:bonding}: contractive downward bonding maps
  $p_{n+1, n}^0$ (restriction) and $p_{n+1, n}^1$ (parent-fibre
  averaging).
\item Proposition~\ref{prop:adjoints}: cylindrical embeddings
  $j_{n, n+1}^0, j_{n, n+1}^1$ as adjoints of the bonding maps,
  with the section relations $p_{n+1, n}^0 \circ j_{n, n+1}^0 =
  \id_{X_n^0}$ (vertex sector) and
  $p_{n+1, n}^1 \circ j_{n, n+1}^1 = \tfrac{1}{2}\,\id_{X_n^1}$
  (edge sector); see also Remark~\ref{rmk:edge-section} for the
  unscaled section $s^1$ in the edge sector.
\item Theorem~\ref{thm:Xcyl-direct-limit}: $X_\cyl^{0, \bullet}
  \cong c_{00}(V_\infty^\bullet; F^0)$ as algebraic direct limit
  (vertex sector).
\item Theorem~\ref{thm:Xinv-Hilbert}: $X_\inv^{0, \bullet} \cong
  \ell^2(V_\infty^\bullet; F^0)$ as separable Hilbert space
  (vertex sector).
\item Proposition~\ref{prop:Xcyl-dense}: $X_\cyl^{0, \bullet}$
  dense in $X_\inv^{0, \bullet}$.
\item Theorem~\ref{thm:commute}: refinement-Coxeter (full),
  refinement-$\sigma$ (full), and Coxeter-$\sigma$ (restricted, see
  Remark~\ref{rmk:coxeter-sigma-restricted}) intertwining
  identities.
\end{enumerate}

\subsection{Recommended citation phrasing}

``By P3 \cite{RefinementSpaces}, the cascade vertex substrate is
$\ell^2(V_\infty^\bullet; F^0)$ as the projective inverse limit
of $(X_n^{0, \bullet}, p_{n+1, n}^0)_n$ with $p_{n+1, n}^0$
restriction (Theorems~\ref{thm:bonding}, \ref{thm:Xinv-Hilbert});
the cylindrical direct limit $X_\cyl^{0, \bullet} \cong
c_{00}(V_\infty^\bullet; F^0)$ embeds densely in $X_\inv^{0,
\bullet}$ (Theorem~\ref{thm:Xcyl-direct-limit},
Proposition~\ref{prop:Xcyl-dense}); the refinement-Coxeter and
refinement-$\sigma$ intertwinings hold without restriction, and
the Coxeter-$\sigma$ intertwining holds in the restricted sense of
Theorem~\ref{thm:commute}(4) and
Remark~\ref{rmk:coxeter-sigma-restricted} (full $W(H_4)$-$\sigma$
commutation on the $V_{H_4}$ sector fails as a structural matter,
because a non-trivial $W(H_4)$ reflection does not preserve the
$\sigma$-fixed $\Q$-form $V_{H_4, \Q} \otimes 1$, so the hypothesis
of \cite[Theorem~\texttt{thm:sigma-commutes}]{SigmaRationality} is not satisfied on that
summand).''

\subsection{What this paper does not provide}

No closure-functional / dynamics theorem (P4); no metric, gauge,
spectral, Hecke, $\sigma$-cohomology, or computation
correspondence (P5--P13); no continuum-limit identification of
$X_\inv$ with any classical $L^2$ space; no physical interpretation.

\appendix

%=====================================================================
\section{Notation Audit}
\label{app:notation}
%=====================================================================

\begin{center}
\begin{tabular}{lll}
\hline
Symbol & Meaning & Defined \\
\hline
$\bullet$ & branch decoration ($H_4$ or $D_4$) & Def.~\ref{def:base-graphs} \\
$V_{600}, V_{24}$ & base vertex sets (P2) & Def.~\ref{def:base-graphs} \\
$G_n^\bullet$ & graph at refinement level $n$ & Def.~\ref{def:refinement} \\
$V(G_n^\bullet)$ & vertex set, with $V_n \subset V_{n+1}$ & Def.~\ref{def:refinement} \\
$M_n^\bullet, C_n^\bullet$ & midpoints and centroids at level $n$ & Def.~\ref{def:refinement} \\
$\pi_{n+1, n}^E$ & canonical edge-parent map & Def.~\ref{def:edge-parent} \\
$\mu_n^\bullet, \nu_n^\bullet$ & vertex / edge counting measures & Def.~\ref{def:counting-measures} \\
$F^0$ & vertex fibre $V_{H_4}\oplus V_{D_4}\oplus V_{16}\oplus V_8$ & Notation~\ref{not:fibres} \\
$F^1$ & edge fibre $\Im(\OO)$ (fixed) & Notation~\ref{not:fibres} \\
$X_n^0, X_n^1$ & finite-level $\ell^2$ spaces & Defs.~\ref{def:Xn-0}, \ref{def:Xn-1} \\
$p_{n+1, n}^0$ & vertex bonding (restriction) & Def.~\ref{def:p0} \\
$p_{n+1, n}^1$ & edge bonding (parent-fibre averaging) & Def.~\ref{def:p1} \\
$j_{n, n+1}^0$ & vertex embedding (zero-extension) & Def.~\ref{def:j} \\
$j_{n, n+1}^1$ & edge embedding (parent-replication) & Def.~\ref{def:j} \\
$V_\infty^\bullet$ & $\bigcup_n V(G_n^\bullet)$ & Thm.~\ref{thm:Xcyl-direct-limit} \\
$X_\cyl^{0, \bullet}$ & $\cong c_{00}(V_\infty^\bullet; F^0)$ & Def.~\ref{def:Xcyl-vertex}, Thm.~\ref{thm:Xcyl-direct-limit} \\
$X_\inv^{0, \bullet}$ & $\cong \ell^2(V_\infty^\bullet; F^0)$ & Def.~\ref{def:Xinv}, Thm.~\ref{thm:Xinv-Hilbert} \\
$\sigma$ & coefficientwise Galois involution per P1 & Def.~\ref{def:sigma} \\
$W^\bullet$ & Coxeter group ($W(H_4)$ or $W(D_4)$) & Def.~\ref{def:base-graphs} \\
\hline
\end{tabular}
\end{center}

\begin{thebibliography}{99}

\bibitem{SigmaRationality}
\emph{Cascade--Classical Correspondence I: Sigma-Fixed Rationality
over $\Q(\varphi)$},
\texttt{papers/cascade-sigma-rationality/cascade-sigma-rationality.tex}.
(Paper P1.) Pinned results used in this paper, cited by label:
Definition~\texttt{def:sigma} (the Galois involution
$\sigma \colon \Q(\varphi) \to \Q(\varphi)$); coefficientwise
$\sigma_V$ on $V \otimes_\Q \Q(\varphi)$ as defined in
\S\texttt{sec:coefficientwise}; functoriality of the scalar
extension and Theorem~\texttt{thm:sigma-commutes} ($\sigma$ commutes
with $\Q$-linear endomorphisms defined over the $\sigma$-fixed
$\Q$-form, after scalar extension).

\bibitem{AlgebraicSubstrate}
\emph{Cascade--Classical Correspondence II: The Cascade Algebraic
Substrate},
\texttt{papers/cascade-algebraic-substrate/cascade-algebraic-substrate.tex}.
(Paper P2.) Pinned results used in this paper, cited by label:
\texttt{thm:icosian-closure} (icosian model of $V_{H_4}$, with
$\Z[\varphi]$-rational coordinates and $\Q(\varphi)$-rational
matrix-entry inference for the standard $W(H_4)$ reflection
generators); \texttt{thm:pi-H} ($V_{600}$ vertex set and explicit
coordinates); \texttt{thm:shell-class} ($V_{24} \subset V_{600}$
shell decomposition); \texttt{def:Cl-1-3-basis} (the $\Cl(1, 3)$
monomial basis $\{\gamma_A\}$); \texttt{def:octonion-basis}
(standard octonion basis); sector fibre conventions for $V_{16}$
and $V_8$.

\bibitem{RefinementSpaces}
(This paper.) \emph{Cascade--Classical Correspondence III:
Refinement Spaces and Symmetry Actions}.

\bibitem{CascadeClosureDynamics}
\emph{Cascade Closure Dynamics: Finite Closure Functionals,
Gradient Flow, and Refinement Compatibility}, paper P4.
\texttt{papers/cascade-closure-dynamics/cascade-closure-dynamics.tex}.
The downstream paper that puts dynamics on the chain spaces
constructed here.

\bibitem{CascadeRefinementCompat}
\emph{The Cascade Curvature Operator under Refinement:
Schur--Complement Identity and Energy Intertwining}.
\texttt{papers/cascade-refinement-compat/cascade-refinement-compat.tex}.
Discharges the residual-monotonicity obligation~(O3) of
\cite{CascadeMechanism} for the scaled vertex Laplacian in an
abstract scalar pure-midpoint refinement model;
\S1.5 of that paper names the open lift hypotheses (L1) and (L3)
needed to extend its results to the full refinement datum
constructed here.

\bibitem{CascadeMechanism}
\emph{Cascade Closure: A Conditional Compatibility Template for
Geometric State Selection}, paper P-mech.
\texttt{papers/cascade-mechanism/cascade-mechanism.tex}.
Defines the consumer API (O0)--(O3) that
\cite{CascadeRefinementCompat} discharges in the abstract model.

\bibitem{ReedSimonI}
M.~Reed and B.~Simon, \emph{Methods of Modern Mathematical
Physics, Vol.~I: Functional Analysis}, rev.\ ed., Academic Press,
1980. Hilbert spaces and direct/inverse limits: Ch.~II;
$\ell^2$-spaces on countable sets: Ch.~II.4; algebraic direct
limits: Ch.~II.1.4.

\bibitem{Bogachev}
V.~I.~Bogachev, \emph{Measure Theory}, Vol.~I, Springer, 2007.
Conditional expectation and projective systems: Ch.~10.

\bibitem{Humphreys}
J.~E.~Humphreys, \emph{Reflection Groups and Coxeter Groups},
Cambridge Studies in Advanced Mathematics \textbf{29}, Cambridge
University Press, 1990. $W(H_4)$ and $W(D_4)$: Ch.~2.

\bibitem{CoxeterRegularPolytopes}
H.~S.~M.~Coxeter, \emph{Regular Polytopes}, 3rd ed., Dover, 1973.
$24$-cell: \S 22.2; $600$-cell coordinates: \S 22.7;
Schl\"{a}fli refinement: \S 8.

\bibitem{Hatcher}
A.~Hatcher, \emph{Algebraic Topology}, Cambridge University
Press, 2002. Finite graphs and $1$-cochain conventions: Ch.~1.

\bibitem{Lounesto}
P.~Lounesto, \emph{Clifford Algebras and Spinors}, 2nd ed., London
Math.\ Soc.\ Lecture Note Series \textbf{286}, Cambridge
University Press, 2001. Hilbert--Schmidt-style trace forms: Ch.~9.

\end{thebibliography}

\end{document}
